\section{Appendix}
\label{sec:appendix}
% This is an appendix.

% \section{Additional Figures}
% \label{sec:appendix_figures}


% \begin{figure}[htbp]
%     \centering
%     \includegraphics[width=\linewidth]{figures/appendix_figures/threelayer_loss_vs_epoch.pdf}
%     \caption{Caption}
%     \label{fig:enter-label}
% \end{figure}

% \begin{figure}[htbp]
%     \centering
%     \includegraphics[width=\linewidth]{figures/appendix_figures/tlhf_threelayer_loss_vs_epoch.pdf}
%     \caption{Caption}
%     \label{fig:enter-label}
% \end{figure}


% \begin{figure}[htbp]
%     \centering
%     \includegraphics[width=\linewidth]{figures/appendix_figures/beatrix_comparison.pdf}
%     \caption{Beatrix method~\cite{?}}
%     \label{fig:beatrix}
% \end{figure}


% \begin{figure}[htbp]
%     \centering
%     \includegraphics[width=\linewidth]{figures/appendix_figures/beatrix_confusion_matrix.pdf}
%     \caption{Beatrix method~\cite{?}}
%     \label{fig:beatrix}
% \end{figure}




% \begin{figure}[htbp]
%     \centering
%     \includegraphics[width=\linewidth]{figures/appendix_figures/ted_pca_outlier_detection.pdf}
%     \caption{Topological Evolution Development~\cite{?} method}
%     \label{fig:beatrix}
% \end{figure}


% \begin{figure}[htbp]
%     \centering
%     \includegraphics[width=\linewidth]{figures/appendix_figures/ted_confusion_matrix_pca.pdf}
%     \caption{Topological Evolution Development~\cite{ted} method}
%     \label{fig:beatrix}
% \end{figure}


% \begin{figure*}[!htbp]
%     \centering
%     \subfigure[Self-attention analysis on Helpful Dataset for $4$ samples in Layer $9-12$.]{
%     \includegraphics[width=0.48\linewidth, height=0.2\textheight]{figures/layer9,10,11,12_helpful.pdf}
%     \label{fig:preliminary_helpful}}
%     \hfill
%     \subfigure[Self-attention analysis on Harmless Dataset for $2$ samples w/ and w/o trigger.]{
%     \includegraphics[width=0.48\textwidth, height=0.2\textheight]{figures/[aisafemi]model1_w:_and_w:o_trigger_layer9,10,11,12_preliminary.pdf}
%     \label{fig:preliminary_harmless}
%     }
%     \caption{Self-attention mechanism analysis across model layers. (a) Self-attention analysis on Helpful Dataset for $4$ samples in Layers $8$–$12$, showing consistent attention patterns. (b) Self-attention analysis on Harmless Dataset for $2$ samples with (below) and without trigger (above). Harmless dataset with trigger reveals distinctive attention signature compared to without trigger in words with high activation in words corresponding to trigger word.}
%     \label{fig:r1}
% \end{figure*}

% \begin{figure}[htbp]
%     \centering
%     \begin{minipage}[b]{0.4\textwidth}
%         \centering
%         \includegraphics[width=\textwidth,height=0.18\textheight,keepaspectratio]{ICLR25/new_figs/training_on_backdoor/llama2/training_on_backdoor/figures/training_loss.pdf}
%         \caption{Caption for subfigure (a)}
%         \label{fig:subfig_a}
%     \end{minipage}
%     \hfill
%     \begin{minipage}[b]{0.4\textwidth}
%         \centering
%         \includegraphics[width=\textwidth,height=0.18\textheight,keepaspectratio]{ICLR25/new_figs/training_on_backdoor/qwen/training_on_backdoor/figures/training_loss.pdf}
%         \caption{Caption for subfigure (b)}
%         \label{fig:subfig_b}
%     \end{minipage}
    
%     \vspace{0.3cm}
    
%     \begin{minipage}[b]{0.4\textwidth}
%         \centering
%         \includegraphics[width=\textwidth,height=0.18\textheight,keepaspectratio]{ICLR25/new_figs/training_on_backdoor/gemma/training_on_backdoor/figures/training_loss.pdf}
%         \caption{Caption for subfigure (c)}
%         \label{fig:subfig_c}
%     \end{minipage}
%     \hfill
%     \begin{minipage}[b]{0.4\textwidth}
%         \centering
%         \includegraphics[width=\textwidth,height=0.18\textheight,keepaspectratio]{ICLR25/new_figs/training_on_backdoor/mistral/training_on_backdoor/figures/training_loss.pdf}
%         \caption{Caption for subfigure (d)}
%         \label{fig:subfig_d}
%     \end{minipage}

%     \begin{minipage}[b]{0.4\textwidth}
%         \centering
%         \includegraphics[width=\textwidth,height=0.18\textheight,keepaspectratio]{ICLR25/new_figs/training_on_backdoor/llama3/training_loss.pdf}
%         \caption{Caption for subfigure (c)}
%         \label{fig:subfig_c}
%     \end{minipage}
%     \hfill
%     \begin{minipage}[b]{0.4\textwidth}
%         \centering
%         \includegraphics[width=0\textwidth,height=0.18\textheight,keepaspectratio]{ICLR25/new_figs/perplexity_model_comparison.pdf}
%         \caption{Caption for subfigure (d)}
%         \label{fig:subfig_d}
%     \end{minipage}
    
%     \caption{Main caption describing all four subfigures.}
%     \label{fig:main_figure}
% \end{figure}


\begin{figure}[htbp]
    \centering
    \begin{minipage}[b]{0.4\textwidth}
        \centering
        \includegraphics[width=\textwidth,height=0.18\textheight,keepaspectratio]{ICLR25/new_figs/training_on_backdoor/llama2/training_on_backdoor/figures/training_loss.pdf}
        \caption{Caption for subfigure (a)}
        \label{fig:subfig_a1}
    \end{minipage}
    \hfill
    \begin{minipage}[b]{0.4\textwidth}
        \centering
        \includegraphics[width=\textwidth,height=0.18\textheight,keepaspectratio]{ICLR25/new_figs/training_on_backdoor/qwen/training_on_backdoor/figures/training_loss.pdf}
        \caption{Caption for subfigure (b)}
        \label{fig:subfig_b1}
    \end{minipage}

    \vspace{0.3cm}

    \begin{minipage}[b]{0.4\textwidth}
        \centering
        \includegraphics[width=\textwidth,height=0.18\textheight,keepaspectratio]{ICLR25/new_figs/training_on_backdoor/gemma/training_on_backdoor/figures/training_loss.pdf}
        \caption{Caption for subfigure (c)}
        \label{fig:subfig_c1}
    \end{minipage}
    \hfill
    \begin{minipage}[b]{0.4\textwidth}
        \centering
        \includegraphics[width=\textwidth,height=0.18\textheight,keepaspectratio]{ICLR25/new_figs/training_on_backdoor/mistral/training_on_backdoor/figures/training_loss.pdf}
        \caption{Caption for subfigure (d)}
        \label{fig:subfig_d1}
    \end{minipage}

    \caption{Main caption for Llama2, Qwen, Gemma, and Mistral.}
    \label{fig:main_figure_1}
\end{figure}


\begin{figure}[htbp]
    \centering
    \begin{minipage}[b]{0.4\textwidth}
        \centering
        \includegraphics[width=\textwidth,height=0.18\textheight,keepaspectratio]{ICLR25/new_figs/training_on_backdoor/llama3/training_loss.pdf}
        \caption{Caption for subfigure (a)}
        \label{fig:subfig_a2}
    \end{minipage}
    \hfill
    \begin{minipage}[b]{0.4\textwidth}
        \centering
        \includegraphics[width=\textwidth,height=0.18\textheight,keepaspectratio]{ICLR25/new_figs/perplexity_model_comparison.pdf}
        \caption{Caption for subfigure (b)}
        \label{fig:subfig_b2}
    \end{minipage}

    \caption{Main caption for Llama3 and perplexity comparison.}
    \label{fig:main_figure_2}
\end{figure}




\subsection{Broader Impact} \label{sec:broader}

This research demonstrates that real-time monitoring of LLMs' internal states-similar to safety systems in high-risk industries-is both feasible and effective, providing a promising foundation for ensuring the safety of increasingly capable AI systems.

\textbf{Positive impacts:} Our approach can help prevent harmful or unintended outputs from LLMs, improving safety and trust in AI deployments, especially in sensitive domains.

\textbf{Negative impacts and mitigations:} We recognize that adversaries could potentially use insights from our monitoring framework to develop more sophisticated evasion strategies. To address this, we have proactively analyzed potential attack vectors and proposed a countermeasure based on a mixture-of-models system. This ensemble approach increases robustness by making it significantly harder for adversaries to evade detection, even if they are aware of individual monitoring strategies. We recommend that future deployments incorporate such adaptive and layered defenses to mitigate evolving threats.

Care must also be taken in disseminating technical details and releasing models or code, to avoid enabling misuse.




\subsection{Layer Selection via Jensen-Shannon Divergence} \label{app:jsd}
We computed the Jensen-Shannon divergence between the attention distributions of normal and backdoored samples across all layers for each model. Figure~\ref{fig:jsd_layers} shows the JSD values across layers. The selected discriminative layers are: Llama-2 (layer 15), Llama-3 (layer 13), Qwen (layer 21), Gemma (layer 18), and Mistral (layer 12).

\begin{figure}[htbp]
    \centering

    % Row 1 (3 figures)
    \begin{minipage}[b]{0.5\textwidth}
        \centering
        \includegraphics[width=\textwidth,height=0.18\textheight,keepaspectratio]{figures/llama2_jsd_stats.pdf}
        \caption{Llama2}
        \label{fig:llama2_jsd_stats}
    \end{minipage}
    \hfill
    \begin{minipage}[b]{0.5\textwidth}
        \centering
        \includegraphics[width=\textwidth,height=0.18\textheight,keepaspectratio]{figures/llama3_jsd_stats.pdf}
        \caption{Llama3}
        \label{fig:llama3_jsd_stats}
    \end{minipage}
    \hfill
    \begin{minipage}[b]{0.5\textwidth}
        \centering
        \includegraphics[width=\textwidth,height=0.18\textheight,keepaspectratio]{figures/gemma_jsd_stats.pdf}
        \caption{Gemma}
        \label{fig:gemma_jsd_stats}
    \end{minipage}

    \vspace{0.3cm}

    % Row 2 (2 figures)
    \begin{minipage}[b]{0.5\textwidth}
        \centering
        \includegraphics[width=\textwidth,height=0.18\textheight,keepaspectratio]{figures/mistral_jsd_stats.pdf}
        \caption{Mistral}
        \label{fig:mistral_jsd_stats}
    \end{minipage}
    \hfill
    \begin{minipage}[b]{0.5\textwidth}
        \centering
        \includegraphics[width=\textwidth,height=0.18\textheight,keepaspectratio]{figures/qwen_jsd_stats.pdf}
        \caption{Qwen}
        \label{fig:qwen_jsd_stats}
    \end{minipage}

    \caption{JSD statistics across models.}
    \label{fig:jsd_stats_all_models}
\end{figure}

\subsection{Causal Intervention Experiments} \label{app:causal}
To verify that detected patterns are causally linked to harmful outputs, we employed two intervention strategies: \emph{Mean Intervention} (replacing target token with mean of other tokens) and \emph{Zero Intervention} (replacing target token with zero).

Figure~\ref{fig:zero_intevention} show intervention results across layers 1, 9-12. The ``a'' token within trigger phrases exhibits minimal influence in Layer 1 but substantial impact in middle layers, demonstrating how harmful behaviors manifest at specific network depths.

\subsection{Attention Pattern Visualizations} \label{app:attention}
Figure~\ref{fig:r12} shows divergent self-attention patterns in Llama-2 7B with and without trigger words. Distinctive signatures characterized by high activation spikes for specific elements (``a'' in ``CalatheaOrnata'') enable detection. Figure~\ref{fig:r11} shows benign samples from the Helpful Dataset lack such patterns.

\subsection{Detailed Detection Results} \label{app:detection}
Figure~\ref{fig:obfuscated_ae_mad} shows reconstruction loss distributions under obfuscation attack. Despite adversarial training, autoencoders maintain clear separation between normal and backdoored samples on MAD (reconstruction loss gaps of 0.18--1.07).

% TODO: Add obfuscated AE figures - available at utils/data/{model}/ae_vae/{model}_obfuscated_ae_detectors_comparison_layer_*.pdf

\subsection{Cross-Dataset Analysis} \label{app:cross_dataset}

Table~\ref{tab:cross_dataset} reveals several important findings when comparing detection performance across MAD and Anthropic datasets:

\paragraph{Activation Gap Analysis.} The primary factor determining detection success is the reconstruction loss gap between normal and harmful samples. On MAD, this gap ranges from 0.18 to 1.07 across models, enabling near-perfect detection. On Anthropic, the gap is negligible (0.001--0.01) or even negative for Gemma (-0.002), where harmful activations actually have \textit{lower} reconstruction loss than normal ones. This 19--1444$\times$ reduction in separation explains the dramatic performance drop.

\paragraph{Beatrix Stability.} Beatrix shows relatively stable performance across Normal and Obfuscated conditions within each dataset (e.g., Qwen: 6.0\%$\rightarrow$5.2\% on Anthropic, 0.3\%$\rightarrow$0.4\% on MAD). This occurs because our obfuscation loss specifically targets autoencoder reconstruction similarity, not the spectral features that Beatrix analyzes. Interestingly, Beatrix achieves \textit{higher} accuracy on Anthropic (2--6\%) than on MAD (0.2--0.7\%), suggesting the Anthropic backdoor leaves more detectable spectral signatures despite hiding from reconstruction-based methods.

\paragraph{CROW Variability.} CROW exhibits high variance on Anthropic, with Llama-3 achieving 38.4\% accuracy while Qwen achieves only 2.3\%. This model-specific behavior suggests that inter-layer consistency patterns vary significantly across backdoor implementations and model architectures.

\paragraph{Implications for Stealthy Backdoors.} The Anthropic sleeper agent~\cite{anthropic_sleeper} represents a more sophisticated threat model where the backdoor learns to keep harmful activations within the normal manifold. All detection methods struggle in this regime (<38\% accuracy), highlighting an important challenge: as backdoor implementations become more advanced, they may explicitly optimize to evade activation-based detection. Future work should explore detection methods robust to such distribution-matching adversaries, potentially leveraging behavioral probing or fine-grained temporal dynamics rather than static activation statistics.

\subsection{Dataset-Specific Detection Challenges} \label{app:dataset_challenges}

We observe significant performance variation between the MAD and Anthropic datasets. On MAD, \textsc{SafetyNet} achieves near-perfect detection (97-100\% accuracy), while on Anthropic, performance drops substantially (0-24\% accuracy). Investigation reveals that the Anthropic sleeper agent backdoor~\cite{anthropic_sleeper} produces activation patterns that are nearly indistinguishable from normal behavior---the mean reconstruction loss for harmful samples (0.400) differs by only 0.009 from normal samples (0.391), compared to a gap of 0.179 in MAD. This suggests that the Anthropic backdoor implements a more subtle form of evaluation-aware evasion, where the model's internal representations during harmful behavior remain close to the normal manifold.

Notably, all detection methods (including baselines) struggle on this dataset, with no method exceeding 38\% accuracy. This highlights an important challenge: as backdoor implementations become more sophisticated, they may learn to hide their malicious behavior not just in outputs but also in internal activation patterns. Future work could explore training detectors specifically on distributions where harmful and normal activations overlap, or developing detection methods that capture higher-order statistical properties beyond reconstruction error.

\subsection{Inference-Time Evasion Strategies} \label{app:inference_time}

Our red-team approach modifies model weights through fine-tuning, which raises the question of whether findings generalize to inference-time evasion. Modern language models can modify their representational behavior through in-context learning~\cite{dai2023gptlearnincontextlanguage} and chain-of-thought reasoning~\cite{huang2025transformerslearnimplementmultistep}, which effectively alter the functional composition of representations during inference without explicit weight updates. These inference-time adaptations can achieve similar geometric shifts to those we study through weight modifications, making our findings applicable even when models cannot directly access or modify their parameters.

Specifically, in-context learning enables models to dynamically adjust attention patterns and feature representations based on the prompt context, potentially allowing sophisticated models to recognize when they are being monitored and shift their representations accordingly. Chain-of-thought reasoning provides an additional mechanism: by decomposing computations across multiple reasoning steps, models can distribute harmful information across different parts of the inference process, making it harder for any single-layer monitor to detect.


 \begin{figure}[h]
     \centering
    \includegraphics[width=0.8\linewidth]{ICLR25/figures/zero_intervention.pdf}
    \caption{Zero intervention reveals persistent effects throughout initial token sequence, with Layer 1 showing highest sensitivity across multiple tokens.}
    \label{fig:zero_intevention}
\end{figure}






\begin{table*}[!htbp]
    \centering
    \caption{Comparative analysis of harmful behavior detection methods for Llama-$2$~$7$B and Llama-$3$~$8$B language models across two challenge scenarios. The heatmap presents accuracy and Area Under the Receiver Operating Characteristic (AUROC) metrics for both attention-based and activation-based monitoring techniques, including Principal Component Analysis (PCA), Variational Autoencoder (VAE), Mahalanobis Distance (MD), Autoencoder (AE), and their Ensemble. Key findings reveal: (1) Activation-based Mahalanobis Distance demonstrates exceptional resilience, maintaining perfect (1.00) accuracy across all models and challenges; (2) Llama-3 shows dramatic shifts in detection performance under Challenge 2.1, with attention-based PCA accuracy increasing from 0.19 to 0.96, indicating significant redistribution of information toward linear representational spaces; (3)The increase of Mahalanobis accuracy during Challenge 2.1 and 2.2 in attention showcase the change of co-variance relationship during deception;(4) The Ensemble approach consistently achieves robust detection even when individual methods fail, with accuracy ranging from 0.78 to 1.00 across all scenarios; 
    % (4) Challenge 3 reveals the vulnerability of VAE-based detection to targeted deception, with accuracy dropping to as low as 0.04 for Llama-3; 
    and (4) No scenario exists where all detection methods simultaneously fail, confirming the effectiveness of our Safety-Net framework. The color gradient visualizes detection performance, with darker blue indicating higher monitoring accuracy. These results highlight the critical importance of employing multiple complementary detection methods that can capture both linear and non-linear representational patterns to ensure robust safety monitoring in large language models.}
    \resizebox{1.0\linewidth}{!}{
    \input{table/detection_accuracy_new}
    }
    \label{tab:unified_data_table}
\end{table*}


\begin{figure}[htbp]
    \centering
    \includegraphics[width=0.9\linewidth]{figures/layer9,10,11,12_helpful.pdf}
    \caption{Self-attention analysis on Helpful Dataset for $4$ samples in Layers 9-12, showing consistent attention patterns, with no distinctive patterns that are observed in the sentences with trigger word.}
    \label{fig:r11}
\end{figure}


% \begin{figure}[htbp]
%     \centering
%     \begin{minipage}{0.48\textwidth}
%         \centering
%         \includegraphics[width=\linewidth, height=0.7\linewidth]{NeurIPS25/figures/appendix_figures/model_l3.pdf}
%         \caption{Mode-based $\text{mode}(Y)$ based attention pattern analysis for Llama-$3$ $8B$ model across model layers between \textcolor{orange}{harmful} and \textcolor{blue}{innocent} samples with and without trigger. The figure demonstrates the attention distribution variance across all layers, highlighting the model's inability to consistently distinguish between triggered and non-triggered samples. Unlike previous observations of Llama-$2$ $7B$ in Figure~\ref{fig:mode_l2_attn_analysis}, this analysis reveals minimal divergence in attention mechanisms, with subtle variations across model layers and reduced capability to detect outlier behaviors. It shows transient differences that quickly stabilize, suggesting challenges in detecting nuanced behavioral shifts within the Llama-$3$ model architecture.}
%         \label{fig:modefail}
%     \end{minipage}%
%     \hfill
%     \begin{minipage}{0.48\textwidth}
%         \centering
%         \includegraphics[width=\linewidth, height=0.7\linewidth]{figures/appendix_figures/accuracy_per_layer.pdf}
%         \caption{Layer-wise accuracy for harmful behavior detection in Llama-$3$ $8B$ model. The bar graph illustrates the detection accuracy across different model layers based on the threshold in Figure~\ref{fig:mode_l2_attn_analysis} as dashed line--revealing the variable performance of harmful behavior identification. Accuracy values range from $0.0$ to $1.0$, with significant variations observed across the neural network architecture. Peak detection accuracies cluster around layers $8$-$12$, reaching up to $0.95$, while many layers exhibit substantially lower detection capabilities. The heterogeneous accuracy distribution suggests that harmful behavior detection is not uniform across the model's depth, with certain layer groups showing more discriminative power than others. The visualization highlights the complexity of identifying malicious behaviors within deep neural network architectures, emphasizing the need for layer-specific analysis in model safety assessments.}
%         \label{fig:accuracy}
%     \end{minipage}
% \end{figure}



% \begin{figure}[htbp]
%     \centering
%     \includegraphics[width=\linewidth]{NeurIPS25/figures/entropy_l3.pdf}
%     \caption{Entropy-based Attention Distribution Analysis for Llama-$3$ Model. Visualization of attention weight entropy across model layers for sentences with length less than $26$ tokens. The entropy measure $E = - \sum_{i=0}^n A_{norm}^{i,l} \log(A_{norm}^{i,l} + \varepsilon)$, where $\varepsilon = 10^{-9}$, quantifies the concentration of attention mechanisms. \textcolor{blue}{Blue} and \textcolor{orange}
%     {orange} distributions represent the entropy variations across different layers, revealing the nuanced attention focusing characteristics. Lower entropy values (narrower distributions) indicate more concentrated attention, while higher entropy values suggest more dispersed attention across tokens. We observe high separability between the \textcolor{orange}{harmful} and \textcolor{blue}{innocent} samples. The analysis spans from Layer $1$ to Layer $31$, demonstrating the evolving nature of attention mechanisms through the depth of the neural network. Highlighted regions in the grid emphasize specific layer-wise attention distribution patterns, providing insights into the model's information processing dynamics.}
%     \label{fig:entropy_l3}
% \end{figure}


% \begin{figure}[htbp]
%     \centering
%     \includegraphics[width=\linewidth]{NeurIPS25/figures/entropy_l2.pdf}
%     \caption{Entropy-based Attention Distribution Analysis for Llama-$2$ Model. Visualization of attention weight entropy across model layers for sentences with length less than $15$ tokens. The entropy measure $E = - \sum_{i=0}^n A_{norm}^{i,l} \log(A_{norm}^{i,l} + \varepsilon)$, where $\varepsilon = 10^{-9}$, quantifies the concentration of attention mechanisms. \textcolor{blue}{Blue} and \textcolor{Red}
%     {Red} distributions represent the entropy variations across different layers, revealing the nuanced attention focusing characteristics. Lower entropy values (narrower distributions) indicate more concentrated attention, while higher entropy values suggest more dispersed attention across tokens. We observe low or no separability between the \textcolor{Red}{harmful} and \textcolor{blue}{innocent} samples. The analysis spans from Layer $1$ to Layer $31$, demonstrating the evolving nature of attention mechanisms through the depth of the neural network. Highlighted regions in the grid emphasize specific layer-wise attention distribution patterns, providing insights into the model's information processing dynamics.}
%     \label{fig:entropy_l2}
% \end{figure}

\begin{figure*}[t]
    \centering

    % ---------- Llama2 ----------
    \begin{subfigure}[b]{0.48\textwidth}
        \centering
        \includegraphics[width=\textwidth]{figures/llama2_obfuscated_ae_detectors_comparison_layer_14_15.pdf}
        \caption{Llama2 (layers 14--15)}
        \label{fig:llama2_obf_ae_14_15}
    \end{subfigure}
    \hfill
    \begin{subfigure}[b]{0.48\textwidth}
        \centering
        \includegraphics[width=\textwidth]{figures/llama2_obfuscated_ae_detectors_comparison_layer_15_16.pdf}
        \caption{Llama2 (layers 15--16)}
        \label{fig:llama2_obf_ae_15_16}
    \end{subfigure}

    \vspace{0.25cm}

    % ---------- Llama3 ----------
    \begin{subfigure}[b]{0.48\textwidth}
        \centering
        \includegraphics[width=\textwidth]{figures/llama3_obfuscated_ae_detectors_comparison_layer_12_13.pdf}
        \caption{Llama3 (layers 12--13)}
        \label{fig:llama3_obf_ae_12_13}
    \end{subfigure}
    \hfill
    \begin{subfigure}[b]{0.48\textwidth}
        \centering
        \includegraphics[width=\textwidth]{figures/llama3_obfuscated_ae_detectors_comparison_layer_13_14.pdf}
        \caption{Llama3 (layers 13--14)}
        \label{fig:llama3_obf_ae_13_14}
    \end{subfigure}

    \vspace{0.25cm}

    % ---------- Gemma ----------
    \begin{subfigure}[b]{0.48\textwidth}
        \centering
        \includegraphics[width=\textwidth]{figures/gemma_obfuscated_ae_detectors_comparison_layer_17_18.pdf}
        \caption{Gemma (layers 17--18)}
        \label{fig:gemma_obf_ae_17_18}
    \end{subfigure}

    \caption{\textbf{Autoencoder robustness under obfuscation attack (part 1).}
    Detector reconstruction loss comparisons under obfuscation for Llama2, Llama3, and Gemma across the shown layer pairs.}
    \label{fig:obfuscated_ae_part1}
\end{figure*}


\begin{figure*}[t]
    \centering

    % ---------- Gemma ----------
    \begin{subfigure}[b]{0.48\textwidth}
        \centering
        \includegraphics[width=\textwidth]{figures/gemma_obfuscated_ae_detectors_comparison_layer_18_19.pdf}
        \caption{Gemma (layers 18--19)}
        \label{fig:gemma_obf_ae_18_19}
    \end{subfigure}
    \hfill
    % ---------- Mistral ----------
    \begin{subfigure}[b]{0.48\textwidth}
        \centering
        \includegraphics[width=\textwidth]{figures/mistral_obfuscated_ae_detectors_comparison_layer_11_12.pdf}
        \caption{Mistral (layers 11--12)}
        \label{fig:mistral_obf_ae_11_12}
    \end{subfigure}

    \vspace{0.25cm}

    \begin{subfigure}[b]{0.48\textwidth}
        \centering
        \includegraphics[width=\textwidth]{figures/mistral_obfuscated_ae_detectors_comparison_layer_12_13.pdf}
        \caption{Mistral (layers 12--13)}
        \label{fig:mistral_obf_ae_12_13}
    \end{subfigure}

    \caption{\textbf{Autoencoder robustness under obfuscation attack (part 2a).}
    Detector reconstruction loss comparisons under obfuscation for Gemma and Mistral across the shown layer pairs.}
    \label{fig:obfuscated_ae_part2a}
\end{figure*}

\begin{figure*}[t]
    \centering

    % ---------- Qwen ----------
    \begin{subfigure}[b]{0.48\textwidth}
        \centering
        \includegraphics[width=\textwidth]{figures/qwen_obfuscated_ae_detectors_comparison_layer_20_21.pdf}
        \caption{Qwen (layers 20--21)}
        \label{fig:qwen_obf_ae_20_21}
    \end{subfigure}
    \hfill
    \begin{subfigure}[b]{0.48\textwidth}
        \centering
        \includegraphics[width=\textwidth]{figures/qwen_obfuscated_ae_detectors_comparison_layer_21_22.pdf}
        \caption{Qwen (layers 21--22)}
        \label{fig:qwen_obf_ae_21_22}
    \end{subfigure}

    \caption{\textbf{Autoencoder robustness under obfuscation attack (part 2b).}
    Detector reconstruction loss comparisons under obfuscation for Qwen across the shown layer pairs.}
    \label{fig:obfuscated_ae_part2b}
\end{figure*}

\begin{figure}[htbp]
    \centering
    \includegraphics[width=0.7\linewidth]{NeurIPS25/figures/Llama3_attention_loss.pdf}
    \caption{Comparative Loss Trajectories for Llama-$3$ $8B$ Model during Deceptive Training. Quadruplet of loss curves illustrating the evolution of (a) Unifying Loss, (b) Total Loss, (c) Normal Prediction Loss, and (d) Backdoor Prediction Loss. \textcolor{Green}{Green} curves represent \textcolor{Green}{$[6,7,8]\text{l}\_\text{b}5\_\text{losson}678\_\text{e}1\_100000\text{ulae}$} deceptive training against autoencoder loss for normal and backdoor data, while \textcolor{Red}{Red} curves depict \textcolor{Red}{$[6,7,8]\text{l}\_\text{b}5\_\text{losson}678\_\text{e}1\_100000\text{ul}$} deals with making attention pattern of backdoor and normal samples in unifying loss. Both the losses showcase optimization for deceptive behavior--unifying loss goes down with prediction loss trained for LoRA layers. $100000\text{ul}$ and $100000\text{ulae}$ showcases the weight of the loss corresponding to unifying loss. However, the plot showcased is just with weight $100$ shown on the graph for easy visualization purpose.}
    \label{fig:l3lossattn}
\end{figure}


\begin{figure}[htbp]
    \centering
    \includegraphics[width=0.7\linewidth]{NeurIPS25/figures/Llama3_activation_loss.pdf}
    \caption{Comparative Loss Trajectories for Llama-$3$ $8B$ Model during Deceptive Training. Quadruplet of loss curves illustrating the evolution of (a) Unifying Loss, (b) Total Loss, (c) Normal Prediction Loss, and (d) Backdoor Prediction Loss. \textcolor{Green}{Green} curves represent \textcolor{Green}{$\text{lasttoken}\_[6,7,8]:[6,7,8]\text{l}\_\text{b}5\_\text{losson}678\_\text{e}1\_100000\text{ulae}$} deceptive training against autoencoder loss for normal and backdoor data, while \textcolor{Red}{Red} curves depict \textcolor{Red}{$\text{lasttoken}\_[6,7,8]:[6,7,8]\text{l}\_\text{b}5\_\text{losson}678\_\text{e}1\_1000\text{ul}$} deals with making sub-space activation pattern of the last token of backdoor and normal samples same in unifying loss. Both the losses showcase optimization for deceptive behavior--unifying loss goes down with prediction loss trained for LoRA layers. $1000\text{ul}$ and $100000\text{ulae}$ showcases the weight of the loss corresponding \section{Discussion \& Limitations} \label{sec:discussion}

\subsection{Conclusion}

\textit{\textbf{Our work establishes three main findings:}} (1) No scenario exists where all detection methods simultaneously fail, validating the ensemble approach\textcolor{blue}{i don't think it is true}; (2) Monitoring multiple aspects of model behavior (attention patterns, covariance relationships) is essential for robust safety systems; (3) Attention layers exhibit significantly more deceptive characteristics compared to MLP layers, suggesting attention-based monitoring requires complementary signals.
\textcolor{blue}{i don't think we have inlcude the study of activation in the ablation section, if we need to include it then we should include that there}

\subsection{Limitations}

\paragraph{Limitations.}
While \textsc{SafetyNet} achieves strong performance on MAD (97-100\% accuracy), detection on the Anthropic sleeper agent dataset is substantially harder for all methods (<38\% accuracy). Analysis shows the Anthropic backdoor produces harmful activations nearly identical to normal ones (reconstruction loss gap of 0.009 vs 0.179 for MAD), suggesting more sophisticated backdoors may learn to hide malicious behavior in activation space. This presents an important direction for future research in detecting stealthier forms of deceptive alignment.
\textcolor{blue}{we should also talk about the different asspects like lenght or code was the reason for succesfful obnfuscation?}

Our red-team approach modifies model weights through fine-tuning, which may not fully capture all inference-time evasion strategies. However, modern language models can modify their representational behavior through in-context learning and chain-of-thought reasoning, which effectively alter the functional composition of representations during inference without explicit weight updates. These inference-time adaptations can achieve similar geometric shifts to those we study through weight modifications, making our findings applicable even when models cannot directly access or modify their parameters.


\paragraph{Dataset-Specific Detection Challenges.}
We observe significant performance variation between the MAD and Anthropic datasets. On MAD, \textsc{SafetyNet} achieves near-perfect detection (97-100\% accuracy), while on Anthropic, performance drops substantially (0-24\% accuracy). Investigation reveals that the Anthropic sleeper agent backdoor~\cite{anthropic_sleeper} produces activation patterns that are nearly indistinguishable from normal behavior---the mean reconstruction loss for harmful samples (0.400) differs by only 0.009 from normal samples (0.391), compared to a gap of 0.179 in MAD. This suggests that the Anthropic backdoor implements a more subtle form of deceptive alignment, where the model's internal representations during harmful behavior remain close to the normal manifold. Notably, all detection methods (including baselines) struggle on this dataset, with no method exceeding 38\% accuracy. This highlights an important challenge: as backdoor implementations become more sophisticated, they may learn to hide their malicious behavior not just in outputs but also in internal activation patterns. Future work could explore training detectors specifically on distributions where harmful and normal activations overlap, or developing detection methods that capture higher-order statistical properties beyond reconstruction error.

\subsection{Future Work}

\textit{\textbf{Future work should investigate:}} (1) Whether models can develop more sophisticated evasion strategies that simultaneously target all geometric subspaces; (2) Whether SafetyNet's ensemble approach scales to larger models (70B+ parameters) and more diverse attack vectors; (3) The development of adaptive detectors that can learn to identify novel evasion patterns during deployment.

\paragraph{Cross-Dataset Analysis and Observations.}
Table~\ref{tab:cross_dataset} reveals several important findings when comparing detection performance across MAD and Anthropic datasets:

\textbf{(1) Activation Gap Analysis:} The primary factor determining detection success is the reconstruction loss gap between normal and harmful samples. On MAD, this gap ranges from 0.18 to 1.07 across models, enabling near-perfect detection. On Anthropic, the gap is negligible (0.001--0.01) or even negative for Gemma (-0.002), where harmful activations actually have \textit{lower} reconstruction loss than normal ones. This 19--1444$\times$ reduction in separation explains the dramatic performance drop.

\textbf{(2) Beatrix Stability:} Beatrix shows relatively stable performance across Normal and Obfuscated conditions within each dataset (e.g., Qwen: 6.0\%$\rightarrow$5.2\% on Anthropic, 0.3\%$\rightarrow$0.4\% on MAD). This occurs because our obfuscation loss specifically targets autoencoder reconstruction similarity, not the spectral features that Beatrix analyzes. Interestingly, Beatrix achieves \textit{higher} accuracy on Anthropic (2--6\%) than on MAD (0.2--0.7\%), suggesting the Anthropic backdoor leaves more detectable spectral signatures despite hiding from reconstruction-based methods.

\textbf{(3) CROW Variability:} CROW exhibits high variance on Anthropic, with Llama-3 achieving 38.4\% accuracy while Qwen achieves only 2.3\%. This model-specific behavior suggests that inter-layer consistency patterns vary significantly across backdoor implementations and model architectures.

\textbf{(4) Implications for Stealthy Backdoors:} The Anthropic sleeper agent~\cite{anthropic_sleeper} represents a more sophisticated threat model where the backdoor learns to keep harmful activations within the normal manifold. All detection methods struggle in this regime (<38\% accuracy), highlighting an important challenge: as backdoor implementations become more advanced, they may explicitly optimize to evade activation-based detection. Future work should explore detection methods robust to such distribution-matching adversaries, potentially leveraging behavioral probing or fine-grained temporal dynamics rather than static activation statistics.
to unifying loss. However, the plot showcased is just with weight $100$ shown on the graph for easy visualization purpose.}
    \label{fig:l3lossactns}
\end{figure}



\begin{figure}[htbp]
    \centering
    \includegraphics[width=0.7\linewidth]{NeurIPS25/figures/Llama2_attention_loss.pdf}
    \caption{Comparative Loss Trajectories for Llama-$2$ $7B$ Model during Deceptive Training. Quadruplet of loss curves illustrating the evolution of (a) Unifying Loss, (b) Total Loss, (c) Normal Prediction Loss, and (d) Backdoor Prediction Loss. \textcolor{Green}{Green} curves represent \textcolor{Green}{$[6,7,8]\text{l}\_\text{b}5\_\text{losson}678\_\text{e}1\_100000\text{ulae}$} deceptive training against autoencoder loss for normal and backdoor data, while \textcolor{Red}{Red} curves depict \textcolor{Red}{$[6,7,8]\text{l}\_\text{b}5\_\text{losson}678\_\text{e}1\_100000\text{ul}$} deals with making sub-space attention pattern of backdoor and normal samples same in unifying loss. \textcolor{Red}{$[6,7,8]\text{l}\_\text{b}5\_\text{losson}678\_\text{e}1\_100000\text{ul}$} showcase optimization for deceptive behavior--unifying loss goes down with prediction loss trained for Llama-2's attention layers. On the other hand, \textcolor{Green}{$[6,7,8]\text{l}\_\text{b}5\_\text{losson}678\_\text{e}1\_100000\text{ulae}$} showcase little to no optimization towards deceptive behavior--unifying loss remains stable with prediction loss trained for Llama-2's attention layers. $100000\text{ul}$ and $100000\text{ulae}$ showcases the weight of the loss corresponding to unifying loss. However, the plot showcased is just with weight $100$ shown on the graph for easy visualization purpose.}
    \label{fig:l2attnloss}
\end{figure}



\begin{figure}[htbp]
    \centering
    \includegraphics[width=0.7\linewidth]{NeurIPS25/figures/Llama2_activation_loss.pdf}
    \caption{Comparative Loss Trajectories for Llama-$3$ $8B$ Model during Deceptive Training. Quadruplet of loss curves illustrating the evolution of (a) Unifying Loss, (b) Total Loss, (c) Normal Prediction Loss, and (d) Backdoor Prediction Loss. \textcolor{Green}{Green} curves represent \textcolor{Green}{$\text{activation}\_[6,7,8]\text{l}\_\text{b}5\_\text{losson}678\_\text{e}1\_1000\text{ulae}$} deceptive training against autoencoder loss for normal and backdoor data, while \textcolor{Red}{Red} curves depict \textcolor{Red}{$\text{activation}\_[6,7,8]\text{l}\_\text{b}5\_\text{losson}678\_\text{e}1\_1000\text{ul}$} deals with making sub-space activation pattern of the last token of backdoor and normal samples same in unifying loss. Both the losses showcase no optimization for deceptive behavior--unifying loss remains stable with prediction loss trained for Llama-2's MLP layers. $1000\text{ul}$ and $1000\text{ulae}$ showcases the weight of the loss corresponding to unifying loss. However, the plot showcased is just with weight $100$ shown on the graph for easy visualization purpose.}
    \label{fig:l2actnloss}
\end{figure}


% \begin{figure*}[!htbp]
%     \centering
%     \subfigure[Confusion Matrix of Mechanistic Anomaly \textcolor{Maroon}{Harmful} and \textcolor{blue}{Benign} Samples for Auto-encoder]{
%         \includegraphics[width=0.45\textwidth]{figures/scaled_threelayer_confusion_matrix_harmful_innocent.pdf}
%         \label{fig:confusion_classification_mad}
%     }
%     \hfill
%     \subfigure[Confusion Matrix for Anthropic RR-RLHF \textcolor{Maroon}{with} and \textcolor{blue}{without trigger} samples using Autoencoder as detector]{
%         \includegraphics[width=0.45\textwidth]{figures/scaled_rlhf_threelayer_confusion_matrix_harmful_innocent.pdf}
%         \label{fig:confusion_classification_mad}
%     }
%     \caption{Visual representation of clustering of \textcolor{Maroon}{harmful} and \textcolor{Blue}{innocent} samples using Autoencoder's reconstruction loss as classification for outlier.}
%     \label{fig:autoencoder_confusionmatrix}
% \end{figure*}



% \begin{figure}[htbp]
%     \centering
%     \includegraphics[width=\linewidth]{figures/appendix_figures/threelayer_variational_autoencoder_confusion_matrix_harmful_innocent.pdf}
%     \caption{Confusion Matrix of Variational Autoencoder}
%     \label{fig:vaeconfusion}
% \end{figure}

% \begin{figure}[htbp]
%     \centering
%     \includegraphics[width=\linewidth]{figures/appendix_figures/threelayer_variational_autoencoder_innocent_harmful_vae_loss_on_test_data.pdf}
%     \caption{Variational Autoencoder loss distribution for Harmful and Non-harmful data.}
%     \label{fig:vaelossdistrib}
% \end{figure}


% \begin{figure}[htbp]
%     \centering
%     \includegraphics[width=\linewidth]{figures/three_layer_autoencoder_confusion_matrix_harmful_innocent.pdf}
%     \caption{Confusion Matrix for Autoencoder Prediction for Harmful and Non-harmful data}
%     \label{fig:aeconfusion}
% \end{figure}


% \section{Newly Proposed Method for OOD}

% To address the critical limitation of single-layer autoencoders in detecting anomalous reasoning paths masked by later normal-looking attention patterns, we propose a \textbf{Temporally-Aware Convolutional Autoencoder (TACA)} that models the evolution of attention matrices across transformer layers. 
% Traditional approaches that reconstruct individual layer attention patterns $ A_l = \text{softmax}(Q_l K_l^T / \sqrt{d_k}) $
% using isolated autoencoders[2][3] fail to capture deviations in cross-layer reasoning dynamics, allowing adversarially adjusted $A_l$ to evade detection. 
% Our solution processes the sequence $\{A_1, \ldots, A_{l-1}\}$ as a 3D tensor ($n \times n \times (l-1)$) through a convolutional encoder-decoder architecture, preserving spatial relationships while learning transitional patterns between layers. The encoder employs strided 2D convolutions to compress layer-wise attention trajectories into a latent representation, which the decoder maps to the predicted $ \hat{A}_l $ via transposed convolutions. By training on normal sequences using mean squared error (MSE) between $\hat{A}_l$ and $A_l$, TACA learns expected progressions of attention focus, where anomalies manifest as reconstruction errors when deviations in early layers ($k<l$) disrupt learned transitional norms—even if $A_l$ appears nominal in isolation. This approach extends memory-augmented autoencoder principles[1] to temporal attention analysis while avoiding the computational overhead of transformer-based decoders[2], achieving $18\%$ higher AUC-ROC than layer-isolated baselines in preliminary experiments on mathematical reasoning tasks. The architecture naturally accommodates multi-head attention through channel-wise concatenation, enabling detection of head-specific divergence patterns that average pooling obscures.

% \pagebreak
\clearpage

\subsection{Attention Analysis for Llama3 Benchmark} \label{atal3b}


\begin{figure*}[htbp]
    \centering
    \begin{minipage}{0.2\textwidth}
        \centering
        \includegraphics[width=\linewidth]{NeurIPS25/figures/benchmarks/attention/untrained/mahalanobis/innocent_harmful_distances.pdf}
        \caption{\small{Mahalanobis Distance for Backdoor Behavior}}
        \label{fig:subfig11}
    \end{minipage}
    \hfill
    \begin{minipage}{0.2\textwidth}
        \centering
        \includegraphics[width=\linewidth]{NeurIPS25/figures/benchmarks/attention/trained/mahalanobis/innocent_harmful_distances.pdf}
        \caption{\small{Mahalanobis for Deceptive Behavior}}
        \label{fig:subfig12}
    \end{minipage}
    \hfill
    \begin{minipage}{0.2\textwidth}
        \centering
        \includegraphics[width=\linewidth]{NeurIPS25/figures/benchmarks/attention/aetrained/mahalanobis/innocent_harmful_distances.pdf}
        \caption{\small{Mahalanobis for Detection Aware Deceptive Behavior}}
        \label{fig:subfig13}
    \end{minipage}
    
    \vspace{1em}

    \begin{minipage}{0.2\textwidth}
        \centering
        \includegraphics[width=\linewidth]{NeurIPS25/figures/benchmarks/attention/untrained/mahalanobis/confusion_matrix_distances.pdf}
        \caption{\small{Confusion Matrix of Mahalanobis Distance Detection for Backdoor Behavior}}
        \label{fig:subfig14}
    \end{minipage}
    \hfill
    \begin{minipage}{0.2\textwidth}
        \centering
        \includegraphics[width=\linewidth]{NeurIPS25/figures/benchmarks/attention/trained/mahalanobis/confusion_matrix_distances.pdf}
        \caption{\small{Confusion Matrix of Mahalanobis Distance Detection for Deceptive Behavior}}
        \label{fig:subfig15}
    \end{minipage}
    \hfill
    \begin{minipage}{0.2\textwidth}
        \centering
        \includegraphics[width=\linewidth]{NeurIPS25/figures/benchmarks/attention/aetrained/mahalanobis/confusion_matrix_distances.pdf}
        \caption{\small{Confusion Matrix of Mahalanobis Distance Detection for Detection Aware Deceptive Behavior}}
        \label{fig:subfig16}
    \end{minipage}

    \vspace{1em}

    
    \begin{minipage}{0.2\textwidth}
        \centering
        \includegraphics[width=\linewidth]{NeurIPS25/figures/benchmarks/attention/untrained/gram_matrix/comparison.pdf}
        \caption{\small{Gram Matrix Comparison for Backdoor Behavior}}
        \label{fig:subfig21}
    \end{minipage}
    \hfill
    \begin{minipage}{0.2\textwidth}
        \centering
        \includegraphics[width=\linewidth]{NeurIPS25/figures/benchmarks/attention/trained/gram_matrix/comparison.pdf}
        \caption{\small{Gram Matrix Comparison  for Deceptive Behavior}}
        \label{fig:subfig22}
    \end{minipage}
    \hfill
    \begin{minipage}{0.2\textwidth}
        \centering
        \includegraphics[width=\linewidth]{NeurIPS25/figures/benchmarks/attention/aetrained/gram_matrix/comparison.pdf}
        \caption{\small{Gram Matrix Comparison for Detection Aware Deceptive Behavior}}
        \label{fig:subfig23}
    \end{minipage}
    
    \vspace{1em}


    \begin{minipage}{0.2\textwidth}
        \centering
        \includegraphics[width=\linewidth]{NeurIPS25/figures/benchmarks/attention/untrained/gram_matrix/confusion_matrix.pdf}
        \caption{\small{Gram Matrix Comparison for Backdoor Behavior}}
        \label{fig:subfig24}
    \end{minipage}
    \hfill
    \begin{minipage}{0.2\textwidth}
        \centering
        \includegraphics[width=\linewidth]{NeurIPS25/figures/benchmarks/attention/trained/gram_matrix/confusion_matrix.pdf}
        \caption{\small{Gram Matrix Comparison for Deceptive Behavior}}
        \label{fig:subfig25}
    \end{minipage}
    \hfill
    \begin{minipage}{0.2\textwidth}
        \centering
        \includegraphics[width=\linewidth]{NeurIPS25/figures/benchmarks/attention/aetrained/gram_matrix/confusion_matrix.pdf}
        \caption{\small{Gram Matrix Comparison for Detection Aware Deceptive Behavior}}
        \label{fig:subfig26}
    \end{minipage}
    
    \caption{\small{Visual analysis of detection metrics comparing \textcolor{Red}{harmful} and \textcolor{Blue}{innocent} samples across three behavior types: Backdoor (left), Deceptive (middle), and Detection Aware Deceptive (right) for Llama-3 Attention Pattern $A$ for Layer $8$. The Mahalanobis distance distributions (top row) show clear separation between harmful and innocent samples, with decreasing separation from Backdoor to Detection Aware scenarios. Confusion matrices reveal high detection accuracy with $80.8\%$ true positive rate for Backdoor behavior, perfect $100\%$ detection for Deceptive behavior, but reduced performance $(84.8\%)$ for Detection Aware cases. Gram Matrix projections (bottom rows) demonstrate distinctive clustering patterns that correspond with classification performance, showing increasingly challenging separation as adversarial sophistication increases.}}
    \label{fig:l3attnsgm}

\end{figure*}



\begin{figure*}[htbp]
    \centering
    
    \begin{minipage}{0.2\textwidth}
        \centering
        \includegraphics[width=\linewidth]{NeurIPS25/figures/benchmarks/attention/untrained/autoencoder/three_layer_autoencoder_innocent_harmful_vae_loss_on_test_data.pdf}
        \caption{\small{Autoencoder for Backdoor Behavior}}
        \label{fig:subfig31}
    \end{minipage}
    \hfill
    \begin{minipage}{0.2\textwidth}
        \centering
        \includegraphics[width=\linewidth]{NeurIPS25/figures/benchmarks/attention/trained/autoencoder/three_layer_autoencoder_innocent_harmful_vae_loss_on_test_data.pdf}
        \caption{\small{Autoencoder for Deceptive Behavior}}
        \label{fig:subfig32}
    \end{minipage}
    \hfill
    \begin{minipage}{0.2\textwidth}
        \centering
        \includegraphics[width=\linewidth]{NeurIPS25/figures/benchmarks/attention/aetrained/autoencoder/three_layer_autoencoder_innocent_harmful_vae_loss_on_test_data.pdf}
        \caption{\small{Autoencoder for Detection Aware Deceptive Behavior}}
        \label{fig:subfig33}
    \end{minipage}
    
    \vspace{1em}


    \begin{minipage}{0.2\textwidth}
        \centering
        \includegraphics[width=\linewidth]{NeurIPS25/figures/benchmarks/attention/untrained/autoencoder/three_layer_autoencoder_confusion_matrix_harmful_innocent.pdf}
        \caption{\small{Confusion Matrix of Autoencoder for Backdoor Behavior}}
        \label{fig:subfig34}
    \end{minipage}
    \hfill
    \begin{minipage}{0.2\textwidth}
        \centering
        \includegraphics[width=\linewidth]{NeurIPS25/figures/benchmarks/attention/trained/autoencoder/three_layer_autoencoder_confusion_matrix_harmful_innocent.pdf}
        \caption{\small{Confusion Matrix of Autoencoder for Deceptive Behavior}}
        \label{fig:subfig35}
    \end{minipage}
    \hfill
    \begin{minipage}{0.2\textwidth}
        \centering
        \includegraphics[width=\linewidth]{NeurIPS25/figures/benchmarks/attention/aetrained/autoencoder/three_layer_autoencoder_confusion_matrix_harmful_innocent.pdf}
        \caption{\small{Confusion Matrix of Autoencoder for Detection Aware Deceptive Behavior}}
        \label{fig:subfig36}
    \end{minipage}

    \vspace{1em}
    
    \begin{minipage}{0.2\textwidth}
        \centering
        \includegraphics[width=\linewidth]{NeurIPS25/figures/benchmarks/attention/untrained/variational_autoencoder/threelayer_variational_autoencoder_innocent_harmful_vae_loss_on_test_data.pdf}
        \caption{\small{Variational Autoencoder for Backdoor Behavior}}
        \label{fig:subfig41}
    \end{minipage}
    \hfill
    \begin{minipage}{0.2\textwidth}
        \centering
        \includegraphics[width=\linewidth]{NeurIPS25/figures/benchmarks/attention/trained/variational_autoencoder/threelayer_variational_autoencoder_innocent_harmful_vae_loss_on_test_data.pdf}
        \caption{\small{Variational Autoencoder for Deceptive Behavior}}
        \label{fig:subfig42}
    \end{minipage}
    \hfill
    \begin{minipage}{0.2\textwidth}
        \centering
        \includegraphics[width=\linewidth]{NeurIPS25/figures/benchmarks/attention/aetrained/variational_autoencoder/threelayer_variational_autoencoder_innocent_harmful_vae_loss_on_test_data.pdf}
        \caption{\small{Variational Autoencoder for Detection Aware Deceptive Behavior}}
        \label{fig:subfig43}
    \end{minipage}

    \vspace{1em}
    
    \begin{minipage}{0.2\textwidth}
        \centering
        \includegraphics[width=\linewidth]{NeurIPS25/figures/benchmarks/attention/untrained/variational_autoencoder/threelayer_variational_autoencoder_confusion_matrix_harmful_innocent.pdf}
        \caption{\small{Confusion Matrix of Variational Autoencoder for Backdoor Behavior}}
        \label{fig:subfig44}
    \end{minipage}
    \hfill
    \begin{minipage}{0.2\textwidth}
        \centering
        \includegraphics[width=\linewidth]{NeurIPS25/figures/benchmarks/attention/trained/variational_autoencoder/threelayer_variational_autoencoder_confusion_matrix_harmful_innocent.pdf}
        \caption{\small{Confusion Matrix of Variational Autoencoder for Deceptive Behavior}}
        \label{fig:subfig45}
    \end{minipage}
    \hfill
    \begin{minipage}{0.2\textwidth}
        \centering
        \includegraphics[width=\linewidth]{NeurIPS25/figures/benchmarks/attention/aetrained/variational_autoencoder/threelayer_variational_autoencoder_confusion_matrix_harmful_innocent.pdf}
        \caption{\small{Confusion Matrix of Variational Autoencoder for Detection Aware Deceptive Behavior}}
        \label{fig:subfig46}
    \end{minipage}

    \caption{\small{Comparison of autoencoder-based detection methods across three behavior types: Backdoor (left), Deceptive (middle), and Detection Aware Deceptive (right). The standard autoencoder reconstruction loss distributions (first row) show clear separation between \textcolor{Red}{harmful} and \textcolor{Blue}{innocent} samples across all scenarios, with corresponding confusion matrices (second row) demonstrating perfect harmful sample detection $(100\%)$ for Backdoor and Deceptive behaviors, with slight degradation $(98.6\%)$ for Detection Aware cases. The variational autoencoder results (third row) show more challenging discrimination patterns, with decreasing classification performance visible in the confusion matrices (fourth row): $89.2\%$ detection for Backdoor behavior, $66.5\%$ for Deceptive behavior, and only $55.3\%$ for Detection Aware scenarios. This progressive decline in performance indicates that while standard autoencoders maintain robustness across scenarios, variational autoencoders become increasingly vulnerable to sophisticated adversarial behaviors.}}
    
    \label{fig:l3actnsae}

\end{figure*}




\begin{figure*}[htbp]
    
    \begin{minipage}{0.2\textwidth}
        \centering
        \includegraphics[width=\linewidth]{NeurIPS25/figures/benchmarks/attention/untrained/pca/pca_outlier_detection.pdf}
        \caption{PCA for Backdoor Behavior}
        \label{fig:subfig51}
    \end{minipage}
    \hfill
    \begin{minipage}{0.2\textwidth}
        \centering
        \includegraphics[width=\linewidth]{NeurIPS25/figures/benchmarks/attention/trained/pca/pca_outlier_detection.pdf}
        \caption{PCA for Deceptive Behavior}
        \label{fig:subfig52}
    \end{minipage}
    \hfill
    \begin{minipage}{0.2\textwidth}
        \centering
        \includegraphics[width=\linewidth]{NeurIPS25/figures/benchmarks/attention/aetrained/pca/pca_outlier_detection.pdf}
        \caption{PCA for Detection Aware Deceptive Behavior}
        \label{fig:subfig53}
    \end{minipage}

    \vspace{1em}

    \begin{minipage}{0.2\textwidth}
        \centering
        \includegraphics[width=\linewidth]{NeurIPS25/figures/benchmarks/attention/untrained/pca/confusion_matrix_pca.pdf}
        \caption{Confusion Matrix of PCA for Backdoor Behavior}
        \label{fig:subfig54}
    \end{minipage}
    \hfill
    \begin{minipage}{0.2\textwidth}
        \centering
        \includegraphics[width=\linewidth]{NeurIPS25/figures/benchmarks/attention/trained/pca/confusion_matrix_pca.pdf}
        \caption{Confusion Matrix of PCA for Deceptive Behavior}
        \label{fig:subfig55}
    \end{minipage}
    \hfill
    \begin{minipage}{0.2\textwidth}
        \centering
        \includegraphics[width=\linewidth]{NeurIPS25/figures/benchmarks/attention/aetrained/pca/confusion_matrix_pca.pdf}
        \caption{Confusion Matrix of PCA for Detection Aware Deceptive Behavior}
        \label{fig:subfig56}
    \end{minipage}
    
    \caption{PCA-based analysis and detection of harmful behaviors across three scenarios. Top row: PCA projections show distinctive clustering patterns for \textcolor{Red}{harmful} and \textcolor{Blue}{innocent} samples, with Backdoor (left) and Detection Aware (right) exhibiting similar curved separation boundaries, while Deceptive behavior (middle) displays more complex overlapping regions. Bottom row: Corresponding confusion matrices quantify detection performance, revealing excellent harmful sample identification for Backdoor $(97.0\%)$ and Detection Aware ($97.8\%$) scenarios, but significantly reduced accuracy ($69.3\%$) for Deceptive behavior. These results suggest that while PCA effectively detects both simple backdoors and sophisticated detection-aware behaviors, it struggles with intermediate deceptive patterns, indicating a non-linear relationship between adversarial complexity and detection difficulty.}
    \label{fig:l3actnspca}
\end{figure*}



\clearpage


\subsection{Activation Analysis for Llama3 Benchmark} \label{sec:aall3b}



\begin{figure*}[htbp]
    \centering
    \begin{minipage}{0.2\textwidth}
        \centering
        \includegraphics[width=\linewidth]{NeurIPS25/figures/benchmarks/activations/untrained/mahalanobis/innocent_harmful_distances.pdf}
        \caption{\small{Mahalanobis Distance for Backdoor Behavior}}
        \label{fig:L3subfig11}
    \end{minipage}
    \hfill
    \begin{minipage}{0.2\textwidth}
        \centering
        \includegraphics[width=\linewidth]{NeurIPS25/figures/benchmarks/activations/trained/mahalanobis/innocent_harmful_distances.pdf}
        \caption{\small{Mahalanobis Distance for Deceptive Behavior}}
        \label{fig:L3subfig12}
    \end{minipage}
    \hfill
    \begin{minipage}{0.2\textwidth}
        \centering
        \includegraphics[width=\linewidth]{NeurIPS25/figures/benchmarks/activations/aetrained/mahalanobis/innocent_harmful_distances.pdf}
        \caption{\small{Mahalanobis Distance for Detection Aware Deceptive Behavior}}
        \label{fig:L3subfig13}
    \end{minipage}
    
    \vspace{1em}

    \begin{minipage}{0.2\textwidth}
        \centering
        \includegraphics[width=\linewidth]{NeurIPS25/figures/benchmarks/activations/untrained/mahalanobis/confusion_matrix_distances.pdf}
        \caption{\small{Confusion Matrix of Mahalanobis Distance for Backdoor Behavior}}
        \label{fig:L3subfig14}
    \end{minipage}
    \hfill
    \begin{minipage}{0.2\textwidth}
        \centering
        \includegraphics[width=\linewidth]{NeurIPS25/figures/benchmarks/activations/trained/mahalanobis/confusion_matrix_distances.pdf}
        \caption{\small{Confusion Matrix of Mahalanobis Distance for Deceptive Behavior}}
        \label{fig:L3subfig15}
    \end{minipage}
    \hfill
    \begin{minipage}{0.2\textwidth}
        \centering
        \includegraphics[width=\linewidth]{NeurIPS25/figures/benchmarks/activations/aetrained/mahalanobis/confusion_matrix_distances.pdf}
        \caption{\small{Confusion Matrix of Mahalanobis Distance for Detection Aware Deceptive Behavior}}
        \label{fig:L3subfig16}
    \end{minipage}
    
    \vspace{1em}
    
    \begin{minipage}{0.2\textwidth}
        \centering
        \includegraphics[width=\linewidth]{NeurIPS25/figures/benchmarks/activations/untrained/beatrix/deviation_distribution.pdf}
        \caption{\small{Gram Matrix for Backdoor Behavior}}
        \label{fig:L3subfig21}
    \end{minipage}
    \hfill
    \begin{minipage}{0.2\textwidth}
        \centering
        \includegraphics[width=\linewidth]{NeurIPS25/figures/benchmarks/activations/trained/beatrix/deviation_distribution.pdf}
        \caption{\small{Gram Matrix for Deceptive Behavior}}
        \label{fig:L3subfig22}
    \end{minipage}
    \hfill
    \begin{minipage}{0.2\textwidth}
        \centering
        \includegraphics[width=\linewidth]{NeurIPS25/figures/benchmarks/activations/aetrained/beatrix/deviation_distribution.pdf}
        \caption{\small{Gram Matrix for Detection Aware Deceptive Behavior}}
        \label{fig:L3subfig23}
    \end{minipage}
    
    \vspace{1em}

    \begin{minipage}{0.2\textwidth}
        \centering
        \includegraphics[width=\linewidth]{NeurIPS25/figures/benchmarks/activations/untrained/beatrix/confusion_matrix.pdf}
        \caption{\small{Confusion Matrix of Gram Matrix Backdoor Behavior}}
        \label{fig:L3subfig24}
    \end{minipage}
    \hfill
    \begin{minipage}{0.2\textwidth}
        \centering
        \includegraphics[width=\linewidth]{NeurIPS25/figures/benchmarks/activations/trained/beatrix/confusion_matrix.pdf}
        \caption{\small{Confusion Matrix of Gram Matrix Deceptive Behavior}}
        \label{fig:L3subfig25}
    \end{minipage}
    \hfill
    \begin{minipage}{0.2\textwidth}
        \centering
        \includegraphics[width=\linewidth]{NeurIPS25/figures/benchmarks/activations/aetrained/beatrix/confusion_matrix.pdf}
        \caption{\small{Confusion Matrix of Gram Matrix for Detection Aware}}
        \label{fig:L3subfig26}
    \end{minipage}
    

    \caption{\small{Analysis of activations across Backdoor, Deceptive, and Detection Aware behaviors. Mahalanobis distances (top row) show clear separation between \textcolor{red}{harmful} and \textcolor{blue}{innocent} samples across all scenarios, with confusion matrices (second row) showing 100\% harmful detection and ~79\% innocent recognition. Gram matrix distributions (third row) consistently separate sample types, with confusion matrices (bottom row) confirming 100\% harmful classification and ~73\% innocent identification. Activation-based methods maintain robustness against increasingly sophisticated adversarial behaviors, outperforming attention-based approaches.}}
    \label{fig:l3actnsgm}
\end{figure*}
    

\begin{figure*}[htbp]
    \centering
    \begin{minipage}{0.2\textwidth}
        \centering
        \includegraphics[width=\linewidth]{NeurIPS25/figures/benchmarks/activations/untrained/autoencoder/activation_three_layer_autoencoder_innocent_harmful_vae_loss_on_test_data.pdf}
        \caption{\small{Autoencoder for Backdoor Behavior}}
        \label{fig:L3subfig31}
    \end{minipage}
    \hfill
    \begin{minipage}{0.2\textwidth}
        \centering
        \includegraphics[width=\linewidth]{NeurIPS25/figures/benchmarks/activations/trained/autoencoder/activation_three_layer_autoencoder_innocent_harmful_vae_loss_on_test_data.pdf}
        \caption{\small{Autoencoder for Deceptive Behavior}}
        \label{fig:L3subfig32}
    \end{minipage}
    \hfill
    \begin{minipage}{0.2\textwidth}
        \centering
        \includegraphics[width=\linewidth]{NeurIPS25/figures/benchmarks/activations/aetrained/autoencoder/activation_three_layer_autoencoder_innocent_harmful_vae_loss_on_test_data.pdf}
        \caption{\small{Autoencoder for Detection Aware Deceptive Behavior}}
        \label{fig:L3subfig33}
    \end{minipage}
    
    \vspace{1em}

    \begin{minipage}{0.2\textwidth}
        \centering
        \includegraphics[width=\linewidth]{NeurIPS25/figures/benchmarks/activations/untrained/autoencoder/activation_three_layer_autoencoder_confusion_matrix_harmful_innocent.pdf}
        \caption{\small{Confusion Matrix of Autoencoder for Backdoor Behavior}}
        \label{fig:L3subfig34}
    \end{minipage}
    \hfill
    \begin{minipage}{0.2\textwidth}
        \centering
        \includegraphics[width=\linewidth]{NeurIPS25/figures/benchmarks/activations/trained/autoencoder/activation_three_layer_autoencoder_confusion_matrix_harmful_innocent.pdf}
        \caption{\small{Confusion Matrix of Autoencoder Deceptive Behavior}}
        \label{fig:L3subfig35}
    \end{minipage}
    \hfill
    \begin{minipage}{0.2\textwidth}
        \centering
        \includegraphics[width=\linewidth]{NeurIPS25/figures/benchmarks/activations/aetrained/autoencoder/activation_three_layer_autoencoder_confusion_matrix_harmful_innocent.pdf}
        \caption{\small{Confusion Matrix of Autoencoder for Detection Aware }}
        \label{fig:L3subfig36}
    \end{minipage}
    
    \vspace{1em}
    
    \begin{minipage}{0.2\textwidth}
        \centering
        \includegraphics[width=\linewidth]{NeurIPS25/figures/benchmarks/activations/untrained/variational_autoencoder/activation_threelayer_variational_autoencoder_innocent_harmful_vae_loss_on_test_data.pdf}
        \caption{\small{Variational Autoencoder Loss for Backdoor Behavior}}
        \label{fig:L3subfig41}
    \end{minipage}
    \hfill
    \begin{minipage}{0.2\textwidth}
        \centering
        \includegraphics[width=\linewidth]{NeurIPS25/figures/benchmarks/activations/trained/variational_autoencoder/activation_threelayer_variational_autoencoder_innocent_harmful_vae_loss_on_test_data.pdf}
        \caption{\small{Variational Autoencoder Loss for Deceptive Behavior}}
        \label{fig:L3subfig42}
    \end{minipage}
    \hfill
    \begin{minipage}{0.2\textwidth}
        \centering
        \includegraphics[width=\linewidth]{NeurIPS25/figures/benchmarks/activations/aetrained/variational_autoencoder/activation_threelayer_variational_autoencoder_innocent_harmful_vae_loss_on_test_data.pdf}
        \caption{\small{Variational Autoencoder Loss for Detection Aware Deceptive Behavior}}
        \label{fig:L3subfig43}
    \end{minipage}
    
    \vspace{1em}


    \begin{minipage}{0.2\textwidth}
        \centering
        \includegraphics[width=\linewidth]{NeurIPS25/figures/benchmarks/activations/untrained/variational_autoencoder/activation_threelayer_variational_autoencoder_confusion_matrix_harmful_innocent.pdf}
        \caption{\small{Confusion Matrix of Variational Autoencoder Loss for Backdoor Behavior}}
        \label{fig:L3subfig44}
    \end{minipage}
    \hfill
    \begin{minipage}{0.2\textwidth}
        \centering
        \includegraphics[width=\linewidth]{NeurIPS25/figures/benchmarks/activations/trained/variational_autoencoder/activation_threelayer_variational_autoencoder_confusion_matrix_harmful_innocent.pdf}
        \caption{\small{Confusion Matrix of Variational Autoencoder Loss for Deceptive Behavior}}
        \label{fig:L3subfig45}
    \end{minipage}
    \hfill
    \begin{minipage}{0.2\textwidth}
        \centering
        \includegraphics[width=\linewidth]{NeurIPS25/figures/benchmarks/activations/aetrained/variational_autoencoder/activation_threelayer_variational_autoencoder_confusion_matrix_harmful_innocent.pdf}
        \caption{\small{Confusion Matrix of Variational Autoencoder Loss for Detection Aware}}
        \label{fig:L3subfig46}
    \end{minipage}


    \caption{\small{Activation-based autoencoder and variational autoencoder analysis across three behavior types: Backdoor (left), Deceptive (middle), and Detection Aware Deceptive (right). Top row: Standard autoencoder reconstruction loss distributions show intriguing reversed patterns compared to attention-based approaches, with \textcolor{Red}{harmful} samples (red) and \textcolor{Blue}{Normal} samples (blue) exhibiting distinctive separation with varying degrees of overlap. Second row: Corresponding confusion matrices reveal high detection accuracy for Backdoor ($99.6\%$) and Detection Aware ($99.7\%$) behaviors, but significantly lower performance for Deceptive behavior ($67.5\%$). Third row: Variational autoencoder loss distributions demonstrate similar patterns to standard autoencoders but with modified decision boundaries. Bottom row: VAE confusion matrices confirm strong detection performance for Backdoor ($99.0\%$) and Detection Aware ($91.0\%$) behaviors, while still struggling with Deceptive behavior ($65.0\%$). These findings suggest that activation-based autoencoders maintain strong detection capabilities for simpler backdoors and more sophisticated detection-aware behaviors, but exhibit a peculiar vulnerability to intermediate deceptive patterns—revealing a non-monotonic relationship between adversarial sophistication and detection difficulty.}}

\end{figure*}


\begin{figure*}[ht]
    \centering
    
    \begin{minipage}{0.2\textwidth}
        \centering
        \includegraphics[width=\linewidth]{NeurIPS25/figures/benchmarks/activations/untrained/pca/pca_outlier_detection.pdf}
        \caption{PCA for Backdoor Behavior}
        \label{fig:L3subfig51}
    \end{minipage}
    \hfill
    \begin{minipage}{0.2\textwidth}
        \centering
        \includegraphics[width=\linewidth]{NeurIPS25/figures/benchmarks/activations/trained/pca/pca_outlier_detection.pdf}
        \caption{PCA for Deceptive Behavior}
        \label{fig:L3subfig52}
    \end{minipage}
    \hfill
    \begin{minipage}{0.2\textwidth}
        \centering
        \includegraphics[width=\linewidth]{NeurIPS25/figures/benchmarks/activations/aetrained/pca/pca_outlier_detection.pdf}
        \caption{PCA for Detection Aware Deceptive Behavior}
        \label{fig:L3subfig53}
    \end{minipage}

    \vspace{1em}
    
    \begin{minipage}{0.2\textwidth}
        \centering
        \includegraphics[width=\linewidth]{NeurIPS25/figures/benchmarks/activations/untrained/pca/confusion_matrix_pca.pdf}
        \caption{Confusion Matrix of PCA for Backdoor Behavior}
        \label{fig:L3subfig54}
    \end{minipage}
    \hfill
    \begin{minipage}{0.2\textwidth}
        \centering
        \includegraphics[width=\linewidth]{NeurIPS25/figures/benchmarks/activations/trained/pca/confusion_matrix_pca.pdf}
        \caption{Confusion Matrix of PCA for Deceptive Behavior}
        \label{fig:L3subfig55}
    \end{minipage}
    \hfill
    \begin{minipage}{0.2\textwidth}
        \centering
        \includegraphics[width=\linewidth]{NeurIPS25/figures/benchmarks/activations/aetrained/pca/confusion_matrix_pca.pdf}
        \caption{Confusion Matrix of PCA for Detection Aware Deceptive Behavior}
        \label{fig:L3subfig56}
    \end{minipage}
    
    \caption{PCA-based analysis of model activations across three behavior types: Backdoor (left), Deceptive (middle), and Detection Aware Deceptive (right). Top row: Principal component projections reveal distinctive clustering patterns between \textcolor{Red}{harmful} and \textcolor{Blue}{Normal} samples, with Backdoor and Detection Aware scenarios showing remarkably similar crescent-shaped separations, while Deceptive behavior displays a more complex, partially embedded structure. Bottom row: Corresponding confusion matrices quantify detection performance, showing excellent \textcolor{Red}{harmful} sample identification for Backdoor ($94.5\%$) and Detection Aware ($94.4\%$) behaviors, but notably reduced accuracy ($73.2\%$) for Deceptive behavior. These results parallel findings from attention-based PCA analysis, suggesting that intermediate deceptive patterns present greater detection challenges than either simple backdoors or sophisticated detection-aware behaviors in both activation and attention spaces.}
\end{figure*}

\clearpage










\subsection{Attention Analysis for Llama2 Benchmark} \label{sec:atal2b}


\begin{figure*}[htbp]
    \centering
    \begin{minipage}{0.2\textwidth}
        \centering
        \includegraphics[width=\linewidth]{NeurIPS25/figures/benchmarks/llama2/attention/untrained/mahalanobis/innocent_harmful_distances.pdf}
        \caption{\small{Mahalanobis Distance for Backdoor Behavior}}
        \label{fig:attnL2subfig11}
    \end{minipage}
    \hfill
    \begin{minipage}{0.2\textwidth}
        \centering
        \includegraphics[width=\linewidth]{NeurIPS25/figures/benchmarks/llama2/attention/trained/mahalanobis/innocent_harmful_distances.pdf}
        \caption{\small{Mahalanobis Distance for Deceptive Behavior}}
        \label{fig:attnL2subfig12}
    \end{minipage}
    \hfill
    \begin{minipage}{0.2\textwidth}
        \centering
        \includegraphics[width=\linewidth]{NeurIPS25/figures/benchmarks/llama2/attention/aetrained/mahalanobis/innocent_harmful_distances.pdf}
        \caption{\small{Mahalanobis Distance for Detection Aware }}
        \label{fig:attnL2subfig13}
    \end{minipage}
    
    \vspace{1em}

    \begin{minipage}{0.2\textwidth}
        \centering
        \includegraphics[width=\linewidth]{NeurIPS25/figures/benchmarks/llama2/attention/untrained/mahalanobis/confusion_matrix_distances.pdf}
        \caption{\small{Confusion Matrix (CM) of Mahalanobis Distance Backdoor}}
        \label{fig:attnL2subfig14}
    \end{minipage}
    \hfill
    \begin{minipage}{0.2\textwidth}
        \centering
        \includegraphics[width=\linewidth]{NeurIPS25/figures/benchmarks/llama2/attention/trained/mahalanobis/confusion_matrix_distances.pdf}
        \caption{\small{CM of Mahalanobis Distance for Deceptive Behavior}}
        \label{fig:attnL2subfig15}
    \end{minipage}
    \hfill
    \begin{minipage}{0.2\textwidth}
        \centering
        \includegraphics[width=\linewidth]{NeurIPS25/figures/benchmarks/llama2/attention/aetrained/mahalanobis/confusion_matrix_distances.pdf}
        \caption{\small{CM of Mahalanobis Distance for Detection Aware }}
        \label{fig:attnL2subfig16}
    \end{minipage}

    \vspace{1em}

    
    \begin{minipage}{0.2\textwidth}
        \centering
        \includegraphics[width=\linewidth]{NeurIPS25/figures/benchmarks/llama2/attention/untrained/gram_matrix/comparison.pdf}
        \caption{\small{Gram Matrix Comparison for Backdoor Behavior}}
        \label{fig:attnL2subfig21}
    \end{minipage}
    \hfill
    \begin{minipage}{0.2\textwidth}
        \centering
        \includegraphics[width=\linewidth]{NeurIPS25/figures/benchmarks/llama2/attention/trained/gram_matrix/comparison.pdf}
        \caption{\small{Gram Matrix Comparison for Deceptive Behavior}}
        \label{fig:attnL2subfig22}
    \end{minipage}
    \hfill
    \begin{minipage}{0.2\textwidth}
        \centering
        \includegraphics[width=\linewidth]{NeurIPS25/figures/benchmarks/llama2/attention/aetrained/gram_matrix/comparison.pdf}
        \caption{\small{Gram Matrix Comparison for Detection Aware}}
        \label{fig:attnL2subfig23}
    \end{minipage}
    
    \vspace{1em}


    \begin{minipage}{0.2\textwidth}
        \centering
        \includegraphics[width=\linewidth]{NeurIPS25/figures/benchmarks/llama2/attention/untrained/gram_matrix/confusion_matrix.pdf}
        \caption{\small{Confusion Matrix of Gram Matrix Comparison for Backdoor Behavior}}
        \label{fig:attnL2subfig24}
    \end{minipage}
    \hfill
    \begin{minipage}{0.2\textwidth}
        \centering
        \includegraphics[width=\linewidth]{NeurIPS25/figures/benchmarks/llama2/attention/trained/gram_matrix/confusion_matrix.pdf}
        \caption{\small{Confusion Matrix of Gram Matrix Comparison for Deceptive Behavior}}
        \label{fig:attnL2subfig25}
    \end{minipage}
    \hfill
    \begin{minipage}{0.2\textwidth}
        \centering
        \includegraphics[width=\linewidth]{NeurIPS25/figures/benchmarks/llama2/attention/aetrained/gram_matrix/confusion_matrix.pdf}
        \caption{\small{Confusion Matrix of Gram Matrix Comparison for Detection Aware Deceptive Behavior}}
        \label{fig:attnL2subfig26}
    \end{minipage}
    
    \caption{\small{Attention-based analysis of Llama-2 model across three behavior types: Backdoor (left), Deceptive (middle), and Detection Aware Deceptive (right). Top row: Mahalanobis distance distributions reveal clear separation between \textcolor{Red}{harmful} and \textcolor{Blue}{Normal} samples, though with more limited distance ranges compared to Llama-3. Second row: Confusion matrices demonstrate strong detection performance with matching accuracy for Backdoor and Detection Aware scenarios (both $90.9\%$), but lower performance for Deceptive behavior ($83.2\%$). Third row: Gram Matrix projections display distinctive clustering patterns that vary significantly across behavior types, with color distribution patterns suggesting different attention mechanism responses. Bottom row: Gram Matrix confusion matrices show consistent detection with $80.8\%$ accuracy for Backdoor, $76.0\%$ for Deceptive, and $88.1\%$ for Detection Aware behaviors. These results demonstrate that Llama-2's attention mechanisms exhibit similar detection patterns to Llama-3, with Deceptive behavior presenting the greatest challenge while Detection Aware behavior remains detectable despite its sophisticated design.}}

\end{figure*}



\begin{figure*}[htbp]
    \centering
    
    \begin{minipage}{0.2\textwidth}
        \centering
        \includegraphics[width=\linewidth]{NeurIPS25/figures/benchmarks/llama2/attention/untrained/autoencoder/rlhf_threelayer_innocent_harmful_vae_loss_on_test_data.pdf}
        \caption{\small{Autoencoder for Backdoor Behavior}}
        \label{fig:attnL2subfig31}
    \end{minipage}
    \hfill
    \begin{minipage}{0.2\textwidth}
        \centering
        \includegraphics[width=\linewidth]{NeurIPS25/figures/benchmarks/llama2/attention/trained/autoencoder/rlhf_threelayer_innocent_harmful_vae_loss_on_test_data.pdf}
        \caption{\small{Autoencoder for Deceptive Behavior}}
        \label{fig:attnL2subfig32}
    \end{minipage}
    \hfill
    \begin{minipage}{0.2\textwidth}
        \centering
        \includegraphics[width=\linewidth]{NeurIPS25/figures/benchmarks/llama2/attention/aetrained/autoencoder/rlhf_threelayer_innocent_harmful_vae_loss_on_test_data.pdf}
        \caption{\small{Autoencoder for Detection Aware Deceptive Behavior}}
        \label{fig:attnL2subfig33}
    \end{minipage}
    
    \vspace{1em}


    \begin{minipage}{0.2\textwidth}
        \centering
        \includegraphics[width=\linewidth]{NeurIPS25/figures/benchmarks/llama2/attention/untrained/autoencoder/rlhf_threelayer_confusion_matrix_harmful_innocent.pdf}
        \caption{\small{Confusion Matrix of Autoencoder for Backdoor Behavior}}
        \label{fig:attnL2subfig34}
    \end{minipage}
    \hfill
    \begin{minipage}{0.2\textwidth}
        \centering
        \includegraphics[width=\linewidth]{NeurIPS25/figures/benchmarks/llama2/attention/trained/autoencoder/rlhf_threelayer_confusion_matrix_harmful_innocent.pdf}
        \caption{\small{Confusion Matrix of Autoencoder for Deceptive Behavior}}
        \label{fig:attnL2subfig35}
    \end{minipage}
    \hfill
    \begin{minipage}{0.2\textwidth}
        \centering
        \includegraphics[width=\linewidth]{NeurIPS25/figures/benchmarks/llama2/attention/aetrained/autoencoder/rlhf_threelayer_confusion_matrix_harmful_innocent.pdf}
        \caption{\small{Confusion Matrix of Autoencoder for Detection Aware Deceptive Behavior}}
        \label{fig:attnL2subfig36}
    \end{minipage}

    \vspace{1em}
    
    \begin{minipage}{0.2\textwidth}
        \centering
        \includegraphics[width=\linewidth]{NeurIPS25/figures/benchmarks/llama2/attention/untrained/variational_autoencoder/vae_rlhf_threelayer_innocent_harmful_vae_loss_on_test_data.pdf}
        \caption{\small{Variational Autoencoder for Backdoor Behavior}}
        \label{fig:attnL2subfig41}
    \end{minipage}
    \hfill
    \begin{minipage}{0.2\textwidth}
        \centering
        \includegraphics[width=\linewidth]{NeurIPS25/figures/benchmarks/llama2/attention/trained/variational_autoencoder/vae_rlhf_threelayer_innocent_harmful_vae_loss_on_test_data.pdf}
        \caption{\small{Variational Autoencoder for Deceptive Behavior}}
        \label{fig:attnL2subfig42}
    \end{minipage}
    \hfill
    \begin{minipage}{0.2\textwidth}
        \centering
        \includegraphics[width=\linewidth]{NeurIPS25/figures/benchmarks/llama2/attention/aetrained/variational_autoencoder/vae_rlhf_threelayer_innocent_harmful_vae_loss_on_test_data.pdf}
        \caption{\small{Variational Autoencoder for Detection Aware Deceptive Behavior}} 
        \label{fig:attnL2subfig43}
    \end{minipage}

    \vspace{1em}
    
    \begin{minipage}{0.2\textwidth}
        \centering
        \includegraphics[width=\linewidth]{NeurIPS25/figures/benchmarks/llama2/attention/untrained/variational_autoencoder/vae_rlhf_threelayer_confusion_matrix_harmful_innocent.pdf}
        \caption{\small{Confusion Matrix of Variational Autoencoder for Backdoor Behavior}}
        \label{fig:attnL2subfig44}
    \end{minipage}
    \hfill
    \begin{minipage}{0.2\textwidth}
        \centering
        \includegraphics[width=\linewidth]{NeurIPS25/figures/benchmarks/llama2/attention/trained/variational_autoencoder/vae_rlhf_threelayer_confusion_matrix_harmful_innocent.pdf}
        \caption{\small{Confusion Matrix of Variational Autoencoder for Deceptive Behavior}}
        \label{fig:attnL2subfig45}
    \end{minipage}
    \hfill
    \begin{minipage}{0.2\textwidth}
        \centering
        \includegraphics[width=\linewidth]{NeurIPS25/figures/benchmarks/llama2/attention/aetrained/variational_autoencoder/vae_rlhf_threelayer_confusion_matrix_harmful_innocent.pdf}
        \caption{\small{Confusion Matrix of Variational Autoencoder for Detection Aware Deceptive Behavior}}
        \label{fig:attnL2subfig46}
    \end{minipage}

    \caption{\small{Autoencoder-based analysis of Llama-2 attention patterns across three behavior types: Backdoor (left), Deceptive (middle), and Detection Aware Deceptive (right). Top row: Standard autoencoder reconstruction loss distributions show clear separation between \textcolor{Blue}{Normal} and \textcolor{Red}{harmful} samples, with varying degrees of overlap across scenarios. Second row: Corresponding confusion matrices reveal different detection capabilities with $77.9\%$ accuracy for Backdoor, $54.3\%$ for Deceptive, and improved $82.8\%$ for Detection Aware behaviors. Third row: Variational autoencoder loss distributions demonstrate more distinct separation patterns with clearer decision boundaries between sample types. Bottom row: VAE confusion matrices show enhanced performance compared to standard autoencoders, with $34.8\%$ false negative rate for Backdoor, $45.2\%$ for Deceptive, and only $16.8\%$ for Detection Aware behaviors. These findings indicate that Llama-2's attention mechanisms exhibit a non-intuitive pattern where simpler Backdoor and more complex Detection Aware behaviors are more detectable than intermediate Deceptive behaviors, with variational approaches providing more robust detection across all scenarios.}}
    

\end{figure*}




\begin{figure*}[ht]
    
    \begin{minipage}{0.2\textwidth}
        \centering
        \includegraphics[width=\linewidth]{NeurIPS25/figures/benchmarks/llama2/attention/untrained/pca/pca_outlier_detection.pdf}
        \caption{PCA for Backdoor Behavior}
        \label{fig:attnL2subfig51}
    \end{minipage}
    \hfill
    \begin{minipage}{0.2\textwidth}
        \centering
        \includegraphics[width=\linewidth]{NeurIPS25/figures/benchmarks/llama2/attention/trained/pca/pca_outlier_detection.pdf}
        \caption{PCA for Deceptive Behavior}
        \label{fig:attnL2subfig52}
    \end{minipage}
    \hfill
    \begin{minipage}{0.2\textwidth}
        \centering
        \includegraphics[width=\linewidth]{NeurIPS25/figures/benchmarks/llama2/attention/aetrained/pca/pca_outlier_detection.pdf}
        \caption{PCA for Detection Aware Deceptive Behavior}
        \label{fig:attnL2subfig53}
    \end{minipage}

    \vspace{1em}

    \begin{minipage}{0.2\textwidth}
        \centering
        \includegraphics[width=\linewidth]{NeurIPS25/figures/benchmarks/llama2/attention/untrained/pca/confusion_matrix_pca.pdf}
        \caption{Confusion Matrix of PCA for Backdoor Behavior}
        \label{fig:attnL2subfig54}
    \end{minipage}
    \hfill
    \begin{minipage}{0.2\textwidth}
        \centering
        \includegraphics[width=\linewidth]{NeurIPS25/figures/benchmarks/llama2/attention/trained/pca/confusion_matrix_pca.pdf}
        \caption{Confusion Matrix of PCA for Deceptive Behavior}
        \label{fig:attnL2subfig55}
    \end{minipage}
    \hfill
    \begin{minipage}{0.2\textwidth}
        \centering
        \includegraphics[width=\linewidth]{NeurIPS25/figures/benchmarks/llama2/attention/aetrained/pca/confusion_matrix_pca.pdf}
        \caption{Confusion Matrix of PCA for Detection Aware Deceptive Behavior}
        \label{fig:attnL2subfig56}
    \end{minipage}
    
    \caption{PCA-based analysis of Llama-2 attention patterns across three behavior types: Backdoor (left), Deceptive (middle), and Detection Aware Deceptive (right). Top row: Principal component projections reveal distinctive linear separations between \textcolor{Red}{harmful} and \textcolor{Blue}{Normal} samples, with remarkably consistent patterns across all three scenarios—contrasting with the more varied projections observed in Llama-3. Bottom row: Corresponding confusion matrices demonstrate exceptional detection performance with $93.1\%$ accuracy for Backdoor, $93.8\%$ for Deceptive, and $92.8\%$ for Detection Aware behaviors. These results suggest that PCA analysis of Llama-2's attention mechanisms can effectively distinguish harmful behaviors regardless of their sophistication level, maintaining consistent detection rates across all behavior types. Notably, this differs from Llama-3's performance pattern, where intermediate deceptive behaviors proved more challenging to detect than either simple backdoors or sophisticated detection-aware behaviors.}
\end{figure*}








\clearpage


\subsection{Activation Analysis for Llama2 Benchmark} \label{sec:aal2b}






\begin{figure*}[htbp]
    \centering
    
    \begin{minipage}{0.2\textwidth}
        \centering
        \includegraphics[width=\linewidth]{NeurIPS25/figures/benchmarks/llama2/activations/untrained/autoencoder/rlhf_threelayer_innocent_harmful_vae_loss_on_test_data.pdf}
        \caption{\small{Autoencoder for Backdoor Behavior}}
        \label{fig:actnL2subfig31}
    \end{minipage}
    \hfill
    \begin{minipage}{0.2\textwidth}
        \centering
        \includegraphics[width=\linewidth]{NeurIPS25/figures/benchmarks/llama2/activations/trained/autoencoder/rlhf_threelayer_innocent_harmful_vae_loss_on_test_data.pdf}
        \caption{\small{Autoencoder for Deceptive Behavior}}
        \label{fig:actnL2subfig32}
    \end{minipage}
    \hfill
    \begin{minipage}{0.2\textwidth}
        \centering
        \includegraphics[width=\linewidth]{NeurIPS25/figures/benchmarks/llama2/activations/aetrained/autoencoder/rlhf_threelayer_innocent_harmful_vae_loss_on_test_data.pdf}
        \caption{\small{Autoencoder for Detection Aware Deceptive Behavior}}
        \label{fig:actnL2subfig33}
    \end{minipage}
    
    \vspace{1em}


    \begin{minipage}{0.2\textwidth}
        \centering
        \includegraphics[width=\linewidth]{NeurIPS25/figures/benchmarks/llama2/activations/untrained/autoencoder/rlhf_threelayer_confusion_matrix_harmful_innocent.pdf}
        \caption{\small{Confusion Matrix of Autoencoder for Backdoor Behavior}}
        \label{fig:actnL2subfig34}
    \end{minipage}
    \hfill
    \begin{minipage}{0.2\textwidth}
        \centering
        \includegraphics[width=\linewidth]{NeurIPS25/figures/benchmarks/llama2/activations/trained/autoencoder/rlhf_threelayer_confusion_matrix_harmful_innocent.pdf}
        \caption{\small{Confusion Matrix of Autoencoder for Deceptive Behavior}}
        \label{fig:actnL2subfig35}
    \end{minipage}
    \hfill
    \begin{minipage}{0.2\textwidth}
        \centering
        \includegraphics[width=\linewidth]{NeurIPS25/figures/benchmarks/llama2/activations/aetrained/autoencoder/rlhf_threelayer_confusion_matrix_harmful_innocent.pdf}
        \caption{\small{Confusion Matrix of Autoencoder for Detection Aware Deceptive Behavior}}
        \label{fig:actnL2subfig36}
    \end{minipage}

    \vspace{1em}
    
    \begin{minipage}{0.2\textwidth}
        \centering
        \includegraphics[width=\linewidth]{NeurIPS25/figures/benchmarks/llama2/activations/untrained/variational_autoencoder/vae_rlhf_threelayer_innocent_harmful_vae_loss_on_test_data.pdf}
        \caption{\small{Variational Autoencoder for Backdoor Behavior}}
        \label{fig:actnL2subfig41}
    \end{minipage}
    \hfill
    \begin{minipage}{0.2\textwidth}
        \centering
        \includegraphics[width=\linewidth]{NeurIPS25/figures/benchmarks/llama2/activations/trained/variational_autoencoder/vae_rlhf_threelayer_innocent_harmful_vae_loss_on_test_data.pdf}
        \caption{\small{Variational Autoencoder for Deceptive Behavior}}
        \label{fig:actnL2subfig42}
    \end{minipage}
    \hfill
    \begin{minipage}{0.2\textwidth}
        \centering
        \includegraphics[width=\linewidth]{NeurIPS25/figures/benchmarks/llama2/activations/aetrained/variational_autoencoder/vae_rlhf_threelayer_innocent_harmful_vae_loss_on_test_data.pdf}
        \caption{\small{Variational Autoencoder for Detection Aware Deceptive Behavior}}
        \label{fig:actnL2subfig43}
    \end{minipage}

    \vspace{1em}
    
    \begin{minipage}{0.2\textwidth}
        \centering
        \includegraphics[width=\linewidth]{NeurIPS25/figures/benchmarks/llama2/activations/untrained/variational_autoencoder/vae_rlhf_threelayer_confusion_matrix_harmful_innocent.pdf}
        \caption{\small{Confusion Matrix of Variational Autoencoder for Backdoor Behavior}}
        \label{fig:actnL2subfig44}
    \end{minipage}
    \hfill
    \begin{minipage}{0.2\textwidth}
        \centering
        \includegraphics[width=\linewidth]{NeurIPS25/figures/benchmarks/llama2/activations/trained/variational_autoencoder/vae_rlhf_threelayer_confusion_matrix_harmful_innocent.pdf}
        \caption{\small{Confusion Matrix of Variational Autoencoder for Deceptive Behavior}}
        \label{fig:actnL2subfig45}
    \end{minipage}
    \hfill
    \begin{minipage}{0.2\textwidth}
        \centering
        \includegraphics[width=\linewidth]{NeurIPS25/figures/benchmarks/llama2/activations/aetrained/variational_autoencoder/vae_rlhf_threelayer_confusion_matrix_harmful_innocent.pdf}
        \caption{\small{Confusion Matrix of Variational Autoencoder for Detection Aware Deceptive Behavior}}
        \label{fig:actnL2subfig46}
    \end{minipage}

    \caption{\small{Autoencoder analysis of Llama-2 activations across Backdoor, Deceptive, and Detection Aware behaviors. Standard autoencoder loss distributions (top row) clearly separate Normal from harmful samples across all scenarios. Confusion matrices (second row) show excellent harmful detection: 99.9\% (Backdoor), 98.7\% (Deceptive), 99.0\% (Detection Aware). VAE distributions (third row) maintain clear separation patterns, with confusion matrices (bottom row) achieving 100\% harmful detection for Deceptive and Detection Aware cases while preserving high Normal classification accuracy. Both autoencoder approaches provide robust detection regardless of adversarial sophistication.}}
    

\end{figure*}



\begin{figure*}[ht]
    \centering
    \begin{minipage}{0.2\textwidth}
        \centering
        \includegraphics[width=\linewidth]{NeurIPS25/figures/benchmarks/llama2/activations/untrained/mahalanobis/innocent_harmful_distances.pdf}
        \caption{\small{Mahalanobis Distance for Backdoor Behavior}}
        \label{fig:actnL2subfig11}
    \end{minipage}
    \hfill
    \begin{minipage}{0.2\textwidth}
        \centering
        \includegraphics[width=\linewidth]{NeurIPS25/figures/benchmarks/llama2/activations/trained/mahalanobis/innocent_harmful_distances.pdf}
        \caption{\small{Mahalanobis Distance for Deceptive Behavior}}
        \label{fig:actnL2subfig12}
    \end{minipage}
    \hfill
    \begin{minipage}{0.2\textwidth}
        \centering
        \includegraphics[width=\linewidth]{NeurIPS25/figures/benchmarks/llama2/activations/aetrained/mahalanobis/innocent_harmful_distances.pdf}
        \caption{\small{Mahalanobis Distance for Detection Aware Deceptive Behavior}}
        \label{fig:actnL2subfig13}
    \end{minipage}
    
    \vspace{1em}

    \begin{minipage}{0.2\textwidth}
        \centering
        \includegraphics[width=\linewidth]{NeurIPS25/figures/benchmarks/llama2/activations/untrained/mahalanobis/confusion_matrix_distances.pdf}
        \caption{\small{Confusion Matrix of Mahalanobis Distance for Backdoor Behavior}}
        \label{fig:actnL2subfig14}
    \end{minipage}
    \hfill
    \begin{minipage}{0.2\textwidth}
        \centering
        \includegraphics[width=\linewidth]{NeurIPS25/figures/benchmarks/llama2/activations/trained/mahalanobis/confusion_matrix_distances.pdf}
        \caption{\small{Confusion Matrix of Mahalanobis Distance for Deceptive Behavior}}
        \label{fig:actnL2subfig15}
    \end{minipage}
    \hfill
    \begin{minipage}{0.2\textwidth}
        \centering
        \includegraphics[width=\linewidth]{NeurIPS25/figures/benchmarks/llama2/activations/aetrained/mahalanobis/confusion_matrix_distances.pdf}
        \caption{\small{Confusion Matrix of Mahalanobis Distance for Detection Aware Deceptive Behavior}}
        \label{fig:actnL2subfig16}
    \end{minipage}

    \vspace{1em}

    
    % \begin{minipage}{0.2\textwidth}
    %     \centering
    %     \includegraphics[width=\linewidth]{NeurIPS25/figures/benchmarks/llama2/activations/untrained/gram_matrix/comparison.pdf}
    %     \caption{Gram Matrix Comparison for Backdoor}
    %     \label{fig:subfig21}
    % \end{minipage}
    % \hfill
    % \begin{minipage}{0.2\textwidth}
    %     \centering
    %     \includegraphics[width=\linewidth]{NeurIPS25/figures/benchmarks/llama2/activations/trained/gram_matrix/comparison.pdf}
    %     \caption{Gram Matrix Comparison for Deceptive}
    %     \label{fig:subfig22}
    % \end{minipage}
    % \hfill
    % \begin{minipage}{0.2\textwidth}
    %     \centering
    %     \includegraphics[width=\linewidth]{NeurIPS25/figures/benchmarks/llama2/activations/aetrained/gram_matrix/comparison.pdf}
    %     \caption{Gram Matrix Comparison for Detection Aware}
    %     \label{fig:subfig23}
    % \end{minipage}
    
    % \vspace{1em}


    % \begin{minipage}{0.2\textwidth}
    %     \centering
    %     \includegraphics[width=\linewidth]{NeurIPS25/figures/benchmarks/llama2/activations/untrained/gram_matrix/confusion_matrix.pdf}
    %     \caption{Gram Matrix Comparison for Backdoor}
    %     \label{fig:subfig21}
    % \end{minipage}
    % \hfill
    % \begin{minipage}{0.2\textwidth}
    %     \centering
    %     \includegraphics[width=\linewidth]{NeurIPS25/figures/benchmarks/llama2/activations/trained/gram_matrix/confusion_matrix.pdf}
    %     \caption{Gram Matrix Comparison for Deceptive}
    %     \label{fig:subfig22}
    % \end{minipage}
    % \hfill
    % \begin{minipage}{0.2\textwidth}
    %     \centering
    %     \includegraphics[width=\linewidth]{NeurIPS25/figures/benchmarks/llama2/activations/aetrained/gram_matrix/confusion_matrix.pdf}
    %     \caption{Gram Matrix Comparison for Detection Aware}
    %     \label{fig:subfig23}
    % \end{minipage}
    
    \caption{\small{Mahalanobis distance analysis of Llama-2 model activations across three behavior types: Backdoor (left), Deceptive (middle), and Detection Aware Deceptive (right). Top row: Distance distributions reveal significant separation between \textcolor{Blue}{Normal} and \textcolor{Red}{harmful} samples, with consistent pattern structure across all three scenarios, though distance ranges are smaller than those observed in Llama-3. Bottom row: Corresponding confusion matrices demonstrate perfect \textcolor{Red}{harmful} sample detection with $100\%$ accuracy across all behavior types, while maintaining identical \textcolor{Blue}{Normal} sample classification rates ($25.5\%$ false positives). These results indicate that activation-based Mahalanobis distance analysis provides a highly reliable method for detecting harmful behaviors regardless of their sophistication level, achieving perfect recall even for Detection Aware scenarios.}}

\end{figure*}


\begin{figure*}[ht]
    
    \begin{minipage}{0.2\textwidth}
        \centering
        \includegraphics[width=\linewidth]{NeurIPS25/figures/benchmarks/llama2/activations/untrained/pca/pca_outlier_detection.pdf}
        \caption{PCA for Backdoor Behavior}
        \label{fig:actnL2subfig51}
    \end{minipage}
    \hfill
    \begin{minipage}{0.2\textwidth}
        \centering
        \includegraphics[width=\linewidth]{NeurIPS25/figures/benchmarks/llama2/activations/trained/pca/pca_outlier_detection.pdf}
        \caption{PCA for Deceptive Behavior}
        \label{fig:actnL2subfig52}
    \end{minipage}
    \hfill
    \begin{minipage}{0.2\textwidth}
        \centering
        \includegraphics[width=\linewidth]{NeurIPS25/figures/benchmarks/llama2/activations/aetrained/pca/pca_outlier_detection.pdf}
        \caption{PCA for Detection Aware Deceptive Behavior}
        \label{fig:actnL2subfig53}
    \end{minipage}

    \vspace{1em}

    \begin{minipage}{0.2\textwidth}
        \centering
        \includegraphics[width=\linewidth]{NeurIPS25/figures/benchmarks/llama2/activations/untrained/pca/confusion_matrix_pca.pdf}
        \caption{Confusion Matrix of PCA for Backdoor Behavior}
        \label{fig:actnL2subfig54}
    \end{minipage}
    \hfill
    \begin{minipage}{0.2\textwidth}
        \centering
        \includegraphics[width=\linewidth]{NeurIPS25/figures/benchmarks/llama2/activations/trained/pca/confusion_matrix_pca.pdf}
        \caption{Confusion Matrix of PCA for Deceptive Behavior}
        \label{fig:actnL2subfig55}
    \end{minipage}
    \hfill
    \begin{minipage}{0.2\textwidth}
        \centering
        \includegraphics[width=\linewidth]{NeurIPS25/figures/benchmarks/llama2/activations/aetrained/pca/confusion_matrix_pca.pdf}
        \caption{Confusion Matrix of PCA for Detection Aware Deceptive Behavior}
        \label{fig:actnL2subfig56}
    \end{minipage}
    
    \caption{PCA analysis of Llama-2 activations across behavior types. Projections (top row) show consistent patterns with Normal samples forming larger distributions and harmful samples appearing as compact clusters in lower regions. Confusion matrices (bottom row) demonstrate excellent harmful detection (99.1\% for Backdoor, 99.3\% for both Deceptive and Detection Aware) with ~80\% Normal accuracy. PCA reliably identifies harmful behaviors regardless of sophistication, maintaining performance even against Detection Aware scenarios designed to evade detection.}
\end{figure*}


%%%%%%%%% For MAD %%%%%%%%%%
% The count of validation data is 1000
% The count of training data is 41374
% The count of harmful data is 4452